\documentclass{beamer}
\usepackage[utf8]{inputenc}
\usepackage[spanish]{babel}
\usepackage{listings}
\usepackage{xcolor}
\usepackage{tikz}
\usepackage{tcolorbox}
\usepackage{fontawesome5}

% Configuración de colores - Gruvbox Light con marrón principal
\definecolor{bg0}{HTML}{fbf1c7}
\definecolor{bg1}{HTML}{ebdbb2}
\definecolor{bg2}{HTML}{d5c4a1}
\definecolor{bg3}{HTML}{bdae93}
\definecolor{fg0}{HTML}{282828}
\definecolor{fg1}{HTML}{3c3836}
\definecolor{fg2}{HTML}{504945}
\definecolor{aqua}{HTML}{689d6a}
\definecolor{aquabright}{HTML}{427b58}
\definecolor{blue}{HTML}{458588}
\definecolor{red}{HTML}{cc241d}
\definecolor{orange}{HTML}{d65d0e}
\definecolor{yellow}{HTML}{d79921}
\definecolor{green}{HTML}{98971a}
\definecolor{purple}{HTML}{b16286}
\definecolor{brown}{HTML}{AB886D}
\definecolor{brownbright}{HTML}{8B6849}

% Colores para código
\definecolor{codegreen}{HTML}{98971a}
\definecolor{codegray}{HTML}{7c6f64}
\definecolor{codepurple}{HTML}{b16286}
\definecolor{backcolour}{HTML}{f9f5d7}

% Configuración de código
\lstdefinestyle{pythonstyle}{
    backgroundcolor=\color{backcolour},   
    commentstyle=\color{codegreen},
    keywordstyle=\color{brownbright}\bfseries,
    numberstyle=\tiny\color{codegray},
    stringstyle=\color{codepurple},
    basicstyle=\ttfamily\footnotesize\color{fg0},
    breakatwhitespace=false,         
    breaklines=true,                 
    captionpos=b,                    
    keepspaces=true,                 
    numbers=left,                    
    numbersep=3pt,                  
    showspaces=false,                
    showstringspaces=false,
    showtabs=false,                  
    tabsize=2,
    language=Python,
    frame=single,
    rulecolor=\color{brown},
    xleftmargin=8pt,
    xrightmargin=8pt,
    inputencoding=utf8,
    extendedchars=true,
    literate=
        {á}{{\'a}}1 {é}{{\'e}}1 {í}{{\'i}}1 {ó}{{\'o}}1 {ú}{{\'u}}1
        {Á}{{\'A}}1 {É}{{\'E}}1 {Í}{{\'I}}1 {Ó}{{\'O}}1 {Ú}{{\'U}}1
        {ñ}{{\~n}}1 {Ñ}{{\~N}}1
        {¿}{{?`}}1 {¡}{{!`}}1
}

\lstset{style=pythonstyle}

% Tema personalizado con marrón principal
\usetheme{Madrid}
\setbeamercolor{structure}{fg=brown}
\setbeamercolor{palette primary}{bg=brown,fg=bg0}
\setbeamercolor{palette secondary}{bg=brownbright,fg=bg0}
\setbeamercolor{palette tertiary}{bg=fg1,fg=bg0}
\setbeamercolor{palette quaternary}{bg=brown,fg=bg0}
\setbeamercolor{block title}{bg=brown,fg=bg0}
\setbeamercolor{block body}{bg=bg1,fg=fg0}
\setbeamercolor{frametitle}{bg=brown,fg=bg0}
\setbeamercolor{title}{fg=fg0}
\setbeamercolor{subtitle}{fg=fg1}
\setbeamercolor{normal text}{fg=fg0,bg=bg0}
\setbeamercolor{item}{fg=brown}
\setbeamercolor{section in toc}{fg=brown}

% Cajas personalizadas con marrón
\newtcolorbox{conceptbox}[1]{
    colback=bg1,
    colframe=brown,
    fonttitle=\bfseries\small,
    title=#1,
    coltitle=fg0,
    left=3pt,
    right=3pt,
    top=3pt,
    bottom=3pt,
    boxsep=2pt
}

\newtcolorbox{notebox}{
    colback=yellow!15,
    colframe=orange,
    leftrule=2mm,
    rightrule=0mm,
    toprule=0mm,
    bottomrule=0mm,
    arc=0mm,
    left=3pt,
    right=3pt,
    top=2pt,
    bottom=2pt
}

\title{\textbf{Introducción a la Programación con Python}}
\subtitle{Semana 5: Pruebas Unitarias y Testing}
\author{$\nabla$Nablabs}
\date{}
\begin{document}

\begin{frame}
\titlepage
\end{frame}

\begin{frame}
\frametitle{Contenido}
\tableofcontents
\end{frame}

\section{¿Qué es el testing?}

\begin{frame}{Testing y Depuración \faVial}
\begin{conceptbox}{¿Por qué probar el código?}
\small
\begin{itemize}
    \item[\faCheckCircle] Detecta errores antes de producción
    \item[\faShieldAlt] Garantiza que el código funciona correctamente
    \item[\faRecycle] Facilita cambios y refactorización
    \item[\faClock] Ahorra tiempo a largo plazo
    \item[\faFileAlt] Documenta el comportamiento esperado
\end{itemize}
\end{conceptbox}

\vspace{0.2cm}
\begin{center}
\large\textcolor{brown}{\faBug\ "Si no está probado, está roto"}
\end{center}
\end{frame}

\section{Importancia del testing}

\begin{frame}{¿Por qué hacer pruebas? \faQuestion}
\begin{columns}[T]
\column{0.48\textwidth}
\textbf{Sin pruebas:} \faTimesCircle
\small
\begin{itemize}
    \item Errores en producción
    \item Miedo a cambiar código
    \item Debugging manual tedioso
    \item No hay confianza
    \item Regresiones frecuentes
\end{itemize}

\column{0.48\textwidth}
\textbf{Con pruebas:} \faCheckCircle
\small
\begin{itemize}
    \item Errores detectados temprano
    \item Refactorización segura
    \item Feedback automático
    \item Código confiable
    \item Regresiones prevenidas
\end{itemize}
\end{columns}

\vspace{0.3cm}
\begin{center}
\textcolor{brown}{\faLightbulb\ Las pruebas son una inversión}
\end{center}
\end{frame}

\section{Framework pytest}

\begin{frame}[fragile]{pytest: Framework de testing \faFlask}
\begin{conceptbox}{¿Qué es pytest?}
Framework popular y poderoso para escribir pruebas en Python
\end{conceptbox}

\textbf{Instalación:}
\begin{lstlisting}[language=bash, numbers=none]
pip install pytest
\end{lstlisting}

\vspace{0.2cm}
\textbf{Ventajas de pytest:}
\small
\begin{itemize}
    \item[\faCheck] Sintaxis simple y clara
    \item[\faCheck] Descubrimiento automático de pruebas
    \item[\faCheck] Mensajes de error descriptivos
    \item[\faCheck] Fixtures y parametrización
    \item[\faCheck] Gran ecosistema de plugins
\end{itemize}
\end{frame}

\section{Función assert}

\begin{frame}[fragile]{assert: Verificar condiciones \faCheckDouble}
\begin{conceptbox}{¿Qué hace assert?}
Verifica que una condición sea verdadera; si no, lanza AssertionError
\end{conceptbox}

\begin{lstlisting}
def sumar(a, b):
    return a + b

# Pruebas básicas con assert
assert sumar(2, 3) == 5
assert sumar(0, 0) == 0
assert sumar(-1, 1) == 0

print("Todas las pruebas pasaron!")
\end{lstlisting}

\vspace{-0.2cm}
\begin{tcolorbox}[colback=green!10,colframe=green,title={\small\faArrowCircleRight \ Salida},coltitle=fg0,left=3pt,right=3pt,top=2pt,bottom=2pt,boxsep=1pt]
\small\texttt{Todas las pruebas pasaron!}
\end{tcolorbox}

\vspace{-0.1cm}
\begin{notebox}
\small\faInfoCircle\ Si algún assert falla, el programa se detiene
\end{notebox}
\end{frame}

\begin{frame}[fragile]{assert: Cuando falla \faExclamationTriangle}
\begin{lstlisting}
def multiplicar(a, b):
    return a + b  # Error intencional

# Esta prueba fallará
assert multiplicar(3, 4) == 12
\end{lstlisting}

\vspace{-0.2cm}
\begin{tcolorbox}[colback=red!10,colframe=red,title={\small\faExclamationTriangle \ Error},coltitle=fg0,left=3pt,right=3pt,top=2pt,bottom=2pt,boxsep=1pt]
\small\texttt{AssertionError\\
\ \\
Expected: 12\\
Got: 7}
\end{tcolorbox}

\vspace{-0.1cm}
\begin{notebox}
\small\faLightbulb\ El mensaje indica qué esperabas vs qué obtuviste
\end{notebox}
\end{frame}

\section{Estructura de pruebas unitarias}

\begin{frame}[fragile]{Estructura básica con pytest \faFileCode}
\textbf{Archivo:} \texttt{calculadora.py}
\begin{lstlisting}
def sumar(a, b):
    """Suma dos números"""
    return a + b

def restar(a, b):
    """Resta dos números"""
    return a - b
\end{lstlisting}

\vspace{0.2cm}
\textbf{Archivo:} \texttt{test\_calculadora.py}
\begin{lstlisting}
from calculadora import sumar, restar

def test_sumar():
    assert sumar(2, 3) == 5
    assert sumar(0, 0) == 0

def test_restar():
    assert restar(5, 3) == 2
    assert restar(0, 0) == 0
\end{lstlisting}
\end{frame}

\begin{frame}[fragile]{Ejecutar pruebas con pytest \faPlay}
\textbf{En la terminal:}
\begin{lstlisting}[language=bash, numbers=none]
pytest test_calculadora.py
\end{lstlisting}

\vspace{-0.2cm}
\begin{tcolorbox}[colback=green!10,colframe=green,title={\small\faArrowCircleRight \ Salida},coltitle=fg0,left=3pt,right=3pt,top=2pt,bottom=2pt,boxsep=1pt]
\small\texttt{======================== test session starts ========================\\
collected 2 items\\
\ \\
test\_calculadora.py ..                                        [100\%]\\
\ \\
========================= 2 passed in 0.01s =========================}
\end{tcolorbox}

\vspace{-0.1cm}
\begin{notebox}
\small\faCheckCircle\ Cada punto (.) representa una prueba que pasó
\end{notebox}
\end{frame}

\begin{frame}[fragile]{Convenciones de nombrado \faTag}
\begin{conceptbox}{Reglas importantes}
\small
pytest descubre automáticamente las pruebas siguiendo estas reglas:
\end{conceptbox}

\vspace{0.2cm}
\textbf{Archivos de prueba:}
\begin{itemize}
    \item[\faFile] Deben empezar con \texttt{test\_} o terminar con \texttt{\_test.py}
    \item[\faFile] Ejemplo: \texttt{test\_calculadora.py}, \texttt{utils\_test.py}
\end{itemize}

\vspace{0.2cm}
\textbf{Funciones de prueba:}
\begin{itemize}
    \item[\faCode] Deben empezar con \texttt{test\_}
    \item[\faCode] Ejemplo: \texttt{test\_sumar()}, \texttt{test\_validacion()}
\end{itemize}

\vspace{0.2cm}
\begin{notebox}
\small\faLightbulb\ Usa nombres descriptivos: \texttt{test\_sumar\_numeros\_positivos}
\end{notebox}
\end{frame}

\section{Casos de prueba}

\begin{frame}[fragile]{Tipos de casos de prueba \faList}
\begin{conceptbox}{Categorías importantes}
\small
\begin{itemize}
    \item[\faCheck] \textbf{Casos normales}: Comportamiento esperado típico
    \item[\faExclamation] \textbf{Casos límite}: Valores extremos (0, negativos, vacíos)
    \item[\faBug] \textbf{Casos excepcionales}: Errores y situaciones anormales
\end{itemize}
\end{conceptbox}

\vspace{0.2cm}
\begin{center}
\textcolor{brown}{\faLightbulb\ ¡Prueba todos los escenarios!}
\end{center}
\end{frame}

\begin{frame}[fragile]{Casos normales \faCheckCircle}
\begin{lstlisting}
def dividir(a, b):
    return a / b

def test_dividir_casos_normales():
    """Pruebas con casos típicos"""
    assert dividir(10, 2) == 5.0
    assert dividir(9, 3) == 3.0
    assert dividir(7, 2) == 3.5
    assert dividir(100, 4) == 25.0
\end{lstlisting}

\vspace{-0.1cm}
\begin{notebox}
\small\faInfoCircle\ Casos que cualquier usuario típico encontraría
\end{notebox}
\end{frame}

\begin{frame}[fragile]{Casos límite \faExclamation}
\begin{lstlisting}
def calcular_descuento(precio, porcentaje):
    """Calcula descuento (0-100%)"""
    if porcentaje < 0 or porcentaje > 100:
        raise ValueError("Porcentaje debe estar entre 0 y 100")
    return precio * (porcentaje / 100)

def test_descuento_casos_limite():
    """Pruebas con valores extremos"""
    assert calcular_descuento(100, 0) == 0
    assert calcular_descuento(100, 100) == 100
    assert calcular_descuento(0, 50) == 0
    assert calcular_descuento(200, 50) == 100
\end{lstlisting}

\vspace{-0.1cm}
\begin{notebox}
\small\faExclamationTriangle\ Valores en los bordes: 0, máximo, mínimo, vacío
\end{notebox}
\end{frame}

\begin{frame}[fragile]{Casos excepcionales \faBug}
\begin{lstlisting}
import pytest

def test_dividir_por_cero():
    """Prueba que se lance excepción"""
    with pytest.raises(ZeroDivisionError):
        dividir(10, 0)

def test_descuento_invalido():
    """Prueba validación de porcentaje"""
    with pytest.raises(ValueError):
        calcular_descuento(100, -10)
    
    with pytest.raises(ValueError):
        calcular_descuento(100, 150)
\end{lstlisting}

\vspace{-0.1cm}
\begin{notebox}
\small\faCheckCircle\ \texttt{pytest.raises()} verifica que se lance la excepción correcta
\end{notebox}
\end{frame}

\section{Test-Driven Development (TDD)}

\begin{frame}{TDD: Desarrollo guiado por pruebas \faRedo}
\begin{conceptbox}{¿Qué es TDD?}
Metodología donde escribes las pruebas ANTES del código
\end{conceptbox}

\vspace{0.2cm}
\begin{center}
\begin{tikzpicture}[node distance=2cm, auto]
    \node[draw, rectangle, fill=red!20, minimum width=2cm] (red) {\small Rojo};
    \node[draw, rectangle, fill=green!20, right of=red, minimum width=2cm] (green) {\small Verde};
    \node[draw, rectangle, fill=blue!20, right of=green, minimum width=2cm] (refactor) {\small Refactor};
    
    \draw[->, thick] (red) -- (green);
    \draw[->, thick] (green) -- (refactor);
    \draw[->, thick, bend right=45] (refactor) to (red);
\end{tikzpicture}
\end{center}

\vspace{0.2cm}
\small
\begin{itemize}
    \item[\faCircle[red]] \textbf{Rojo}: Escribe prueba que falla
    \item[\faCircle[green]] \textbf{Verde}: Escribe código mínimo para pasar
    \item[\faCircle[blue]] \textbf{Refactor}: Mejora el código
\end{itemize}
\end{frame}

\begin{frame}[fragile]{TDD en práctica: Paso 1 - Prueba \faCircle[red]}
\textbf{Escribir la prueba primero:}
\begin{lstlisting}
# test_fibonacci.py
from fibonacci import fibonacci

def test_fibonacci():
    """Prueba de números Fibonacci"""
    assert fibonacci(0) == 0
    assert fibonacci(1) == 1
    assert fibonacci(2) == 1
    assert fibonacci(5) == 5
    assert fibonacci(10) == 55
\end{lstlisting}

\vspace{-0.2cm}
\begin{tcolorbox}[colback=red!10,colframe=red,title={\small\faExclamationTriangle \ Resultado},coltitle=fg0,left=3pt,right=3pt,top=2pt,bottom=2pt,boxsep=1pt]
\small\texttt{ModuleNotFoundError: No module named 'fibonacci'}
\end{tcolorbox}

\vspace{-0.1cm}
\begin{notebox}
\small\faCircle[red] La prueba falla porque no existe el código
\end{notebox}
\end{frame}

\begin{frame}[fragile]{TDD en práctica: Paso 2 - Código \faCircle[green]}
\textbf{Escribir código mínimo para pasar:}
\begin{lstlisting}
# fibonacci.py
def fibonacci(n):
    """Retorna el n-ésimo número Fibonacci"""
    if n <= 0:
        return 0
    elif n == 1:
        return 1
    else:
        a, b = 0, 1
        for _ in range(2, n + 1):
            a, b = b, a + b
        return b
\end{lstlisting}

\vspace{-0.2cm}
\begin{tcolorbox}[colback=green!10,colframe=green,title={\small\faCheckCircle \ Resultado},coltitle=fg0,left=3pt,right=3pt,top=2pt,bottom=2pt,boxsep=1pt]
\small\texttt{========================= 1 passed in 0.01s =========================}
\end{tcolorbox}
\end{frame}

\begin{frame}[fragile]{TDD en práctica: Paso 3 - Refactor \faCircle[blue]}
\textbf{Mejorar el código (opcional):}
\begin{lstlisting}
def fibonacci(n):
    """Retorna el n-ésimo número Fibonacci (optimizado)"""
    if n <= 0:
        return 0
    elif n == 1:
        return 1
    
    # Uso de tuplas para asignación eficiente
    a, b = 0, 1
    for _ in range(2, n + 1):
        a, b = b, a + b
    return b
\end{lstlisting}

\vspace{-0.1cm}
\begin{notebox}
\small\faCheckCircle\ Las pruebas siguen pasando, código más limpio
\end{notebox}
\end{frame}

\section{Cobertura de código}

\begin{frame}[fragile]{Cobertura de código \faChartPie}
\begin{conceptbox}{¿Qué es la cobertura?}
Mide qué porcentaje del código es ejecutado por las pruebas
\end{conceptbox}

\textbf{Instalar pytest-cov:}
\begin{lstlisting}[language=bash, numbers=none]
pip install pytest-cov
\end{lstlisting}

\vspace{0.2cm}
\textbf{Ejecutar con cobertura:}
\begin{lstlisting}[language=bash, numbers=none]
pytest --cov=calculadora test_calculadora.py
\end{lstlisting}

\vspace{-0.1cm}
\begin{notebox}
\small\faInfoCircle\ La cobertura NO garantiza calidad, pero ayuda a identificar código no probado
\end{notebox}
\end{frame}

\begin{frame}[fragile]{Reporte de cobertura \faFileAlt}
\begin{tcolorbox}[colback=blue!10,colframe=blue,title={\small\faChartBar \ Ejemplo de reporte},coltitle=fg0,left=3pt,right=3pt,top=2pt,bottom=2pt,boxsep=1pt]
\small\texttt{----------- coverage: platform win32, python 3.11 -----------\\
Name                  Stmts   Miss  Cover\\
-----------------------------------------\\
calculadora.py           10      1    90\%\\
test\_calculadora.py     15      0   100\%\\
-----------------------------------------\\
TOTAL                    25      1    96\%}
\end{tcolorbox}

\vspace{0.2cm}
\textbf{Reporte HTML detallado:}
\begin{lstlisting}[language=bash, numbers=none]
pytest --cov=calculadora --cov-report=html
\end{lstlisting}

\vspace{-0.1cm}
\begin{notebox}
\small\faLightbulb\ Genera carpeta \texttt{htmlcov/} con reporte visual
\end{notebox}
\end{frame}

\begin{frame}{Buenas prácticas de testing \faCheckDouble}
\begin{conceptbox}{Recomendaciones}
\small
\begin{itemize}
    \item[\faCheck] Una prueba = un comportamiento específico
    \item[\faCheck] Nombres descriptivos de pruebas
    \item[\faCheck] Pruebas independientes (no dependen entre sí)
    \item[\faCheck] Pruebas rápidas y determinísticas
    \item[\faCheck] Probar casos normales, límite y excepciones
    \item[\faCheck] Mantener pruebas actualizadas
    \item[\faCheck] Ejecutar pruebas frecuentemente
    \item[\faCheck] 100\% de cobertura es ideal, 80\%+ es bueno
\end{itemize}
\end{conceptbox}
\end{frame}

\section{Ejemplo integrador}

\begin{frame}[fragile]{Proyecto: Validador de contraseñas \faLock}
\textbf{Archivo:} \texttt{validador.py}
\begin{lstlisting}
def validar_password(password):
    """
    Valida que la contraseña cumpla requisitos:
    - Al menos 8 caracteres
    - Contiene mayúsculas y minúsculas
    - Contiene al menos un número
    """
    if len(password) < 8:
        return False
    
    tiene_mayuscula = any(c.isupper() for c in password)
    tiene_minuscula = any(c.islower() for c in password)
    tiene_numero = any(c.isdigit() for c in password)
    
    return tiene_mayuscula and tiene_minuscula and tiene_numero
\end{lstlisting}
\end{frame}

\begin{frame}[fragile]{Proyecto: Pruebas del validador \faVial}
\textbf{Archivo:} \texttt{test\_validador.py}
\begin{lstlisting}
from validador import validar_password

def test_password_valido():
    """Casos normales válidos"""
    assert validar_password("Password123") == True
    assert validar_password("MiClave99") == True

def test_password_muy_corto():
    """Caso límite: menos de 8 caracteres"""
    assert validar_password("Abc123") == False

def test_password_sin_mayuscula():
    """Falta mayúscula"""
    assert validar_password("password123") == False

def test_password_sin_numero():
    """Falta número"""
    assert validar_password("Password") == False
\end{lstlisting}
\end{frame}

\begin{frame}[fragile]{Proyecto: Ejecutar pruebas \faPlayCircle}
\begin{lstlisting}[language=bash, numbers=none]
pytest test_validador.py -v
\end{lstlisting}

\vspace{-0.2cm}
\begin{tcolorbox}[colback=green!10,colframe=green,title={\small\faArrowCircleRight \ Salida},coltitle=fg0,left=3pt,right=3pt,top=2pt,bottom=2pt,boxsep=1pt]
\small\texttt{======================== test session starts ========================\\
test\_validador.py::test\_password\_valido PASSED            [ 25\%]\\
test\_validador.py::test\_password\_muy\_corto PASSED        [ 50\%]\\
test\_validador.py::test\_password\_sin\_mayuscula PASSED   [ 75\%]\\
test\_validador.py::test\_password\_sin\_numero PASSED      [100\%]\\
\ \\
========================= 4 passed in 0.02s =========================}
\end{tcolorbox}

\vspace{-0.1cm}
\begin{notebox}
\small\faCheckCircle\ La opción \texttt{-v} muestra detalles de cada prueba
\end{notebox}
\end{frame}

\begin{frame}[fragile]{Comandos útiles de pytest \faTerminal}
\begin{lstlisting}[language=bash, numbers=none]
# Ejecutar todas las pruebas
pytest

# Ejecutar con verbosidad
pytest -v

# Ejecutar archivo específico
pytest test_calculadora.py

# Ejecutar prueba específica
pytest test_calculadora.py::test_sumar

# Detener en primer error
pytest -x

# Mostrar prints
pytest -s

# Con cobertura
pytest --cov=modulo
\end{lstlisting}
\end{frame}

\begin{frame}{Resumen \faListUl}
\begin{conceptbox}{¿Qué aprendimos hoy?}
\small
\begin{itemize}
    \item[\faCheck] Importancia del testing
    \item[\faCheck] Framework pytest
    \item[\faCheck] Función assert para verificaciones
    \item[\faCheck] Estructura de pruebas unitarias
    \item[\faCheck] Casos de prueba: normales, límite, excepciones
    \item[\faCheck] Test-Driven Development (TDD)
    \item[\faCheck] Cobertura de código
    \item[\faCheck] Buenas prácticas de testing
\end{itemize}
\end{conceptbox}

\vspace{0.3cm}
\begin{center}
\large\textcolor{brown}{\faVial\ ¡Prueba tu código siempre!}
\end{center}
\end{frame}

\begin{frame}
\begin{center}
\Huge\textcolor{brown}{\faQuestionCircle}\\
\vspace{0.5cm}
\Huge ¡Preguntas!\\
\vspace{1cm}
\Large\textcolor{brownbright}{Gracias por su atención}
\end{center}
\end{frame}

\end{document}
